\chapter{What to build, in what order}
\label{sec:roadmap}

The survey supports a fairly confident ordering of implementation work. The
principle behind the ordering: ship the methods with the strongest
evidence-per-line-of-code first, and let each tier's infrastructure carry
the next tier. Every item below names the chapter that justifies it, so
that when the evidence changes, and Chapter~\ref{sec:libraries} is a
reminder of how fast this field's ground shifts, we can trace which
decisions the change touches. Table~\ref{tab:decisions} collects every
verdict in one place; the decision boxes at the ends of
Chapters~\ref{sec:modelfree} through~\ref{sec:workloads} carry the
one-paragraph rationales, and the listed chapter holds the full teaching
and evidence for anyone who wants to challenge a line.

\begin{table}[tp]
\centering
\small
\caption{Every build decision in one place. ``Cost'' is engineering
effort to a first useful version; tiers are this chapter's build order;
the last column is the chapter that teaches the method and holds the
evidence for the verdict. Read a row, then challenge it against its
chapter's decision box.}
\label{tab:decisions}
\begin{tabularx}{\textwidth}{@{}l>{\raggedright\arraybackslash}X>{\raggedright\arraybackslash}Xlll@{}}
\toprule
Capability & What it buys us & When it applies & Cost & Verdict & Ch. \\
\midrule
Sobol + log sampling & honest coverage of every axis & always; also the GP's opening phase & days & T0 & \ref{sec:modelfree} \\
Random-search baseline & the bar every searcher must clear & validation suite & days & T0 & \ref{sec:modelfree} \\
ASHA + median rule & kills bad trials early; the biggest single saving & any workload with learning curves & days & T0 & \ref{sec:multifidelity} \\
TPE (wrap Optuna) & conditional and categorical spaces from few trials & mixed spaces, cheaper trials & days & T0 & \ref{sec:bayesian} \\
GP + qLogEI (BoTorch) & the most information per expensive trial & costly, noisy trials & weeks & T0 & \ref{sec:bayesian} \\
Dimension-scaled priors & reliability past ten dimensions & always on & hours & T0 & \ref{sec:bayesian} \\
qLogNEHVI & noisy Pareto fronts, two to four objectives & Hamiltonian, finance & weeks & T1 & \ref{sec:moo} \\
Chebyshev scalarization & fallback at any objective count; random weights approximate the front when best-seen values proxy the ideal point & five-plus objectives, or cheap engine & days & T1 & \ref{sec:moo} \\
MO-ASHA + veto gates & early stopping for fronts; correctness gates & Hamiltonian, samplers & week & T1 & \ref{sec:moo} \\
$\pi$BO + PriorBand & expert head start that forgets bad advice & every study & 1--2 wks & T1 & \ref{sec:transfer} \\
Warm start from run store & recurring studies start smart & repeated model families & days & T1 & \ref{sec:transfer} \\
EI-per-cost + cooling & budget-aware proposals & unevenly priced trials & days & T1 & \ref{sec:transfer} \\
PB2-Mix, then BG-PBT & schedules and architecture at scale & RL with 8+ parallel workers & months & T2 & \ref{sec:population} \\
DEHB & most consistent small-budget RL tuner & 10--64 full runs & weeks & T3 & \ref{sec:multifidelity} \\
ifBO scheduler & curve-shape-aware budget allocation & once pause/resume is solid & weeks & T3 & \ref{sec:multifidelity} \\
HEBO-style robustness & warping and ensembles on the GP & messy real objectives & weeks & T3 & \ref{sec:bayesian} \\
CARBS cost frontier & tuning while a model scales & scaling-bound projects & later & hold & \ref{sec:transfer} \\
One-shot / zero-cost / grammar NAS & see Section~\ref{sec:computebill}'s verdicts & no current project & --- & skip & \ref{sec:workloads} \\
\bottomrule
\end{tabularx}
\end{table}

\section{Tier 0: wire up what exists (weeks)}

The execution substrate plus existing searchers already covers the sturdy
defaults from Chapters~\ref{sec:modelfree} to~\ref{sec:multifidelity}:
Sobol and log-uniform random sampling (the two free upgrades of
Chapter~\ref{sec:modelfree}), ASHA and the median stopping rule, TPE via
Optuna's sampler for conditional spaces, CMA-ES for cheap continuous
problems, and a single-objective GP searcher with qLogEI (either BoTorch
directly or Optuna's GP sampler), with the dimension-scaled lengthscale
priors of Section~\ref{sec:highdim} switched on from the start, since
they cost one line. None of this is novel; all of it is table stakes,
and the deliverable of this tier is really our own searcher and
scheduler interfaces (Chapter~\ref{sec:architecture}) proven against
known-good implementations, plus the run store that records every trial
with its configuration, seed, fidelity, curve, and cost. A tier is done
when the validation suite below says our wiring matches the reference
implementations, not when the code merely runs.

\section{Tier 1: the differentiators (the core product)}

Four capabilities justify building something of our own, because no single
existing library ships them together on our execution substrate.

Multi-objective, properly. A qLogNEHVI searcher for two to four objectives,
Chebyshev scalarization as the many-objective and cheap fallback (the
division Figure~\ref{fig:chebyshev} argued), MO-ASHA as the
multi-fidelity scheduler, hypervolume-over-budget reporting so that a
study's progress is one honest curve, and Pareto fronts as first-class
study outputs (Chapter~\ref{sec:moo}). This is the piece our Hamiltonian
and finance workloads ask for most directly, and the piece the audit of
Chapter~\ref{sec:libraries} found no executor shipping natively.

Priors and warm starts. The $\pi$BO acquisition weight~\eqref{eq:pibo}
and the PriorBand sampling portfolio over ASHA rungs, both consuming a
prior declared in the search-space schema; plus quantile-based warm
starting from the run store (Chapter~\ref{sec:transfer}). Cheap to
build, immediately felt by every user, and the mechanism's built-in
decay means a wrong prior costs a study its opening moves rather than
its conclusion provided the product schema requires a lower-bounded
mixture or explicitly nonzero support everywhere, so that wrong priors
remain wrong rather than becoming unfalsifiable
(Chapter~\ref{sec:transfer}'s decision box carries this guard).

Cost-aware acquisition. EI-per-cost with cost cooling
\citep{lee2020carbo} on top of the Tier 0 GP searcher. This needs a
\emph{predictive} cost model, because acquisition prices a candidate
before it runs and realized cost cannot do that; what the executor
supplies is the training data for the model, not the price itself. The
specification is deliberately modest: a second GP on
$\log(\text{wall-clock seconds})$, fit to the executor's measured
per-trial costs from the run store, with the same kernel family as the
performance surrogate; log space because cost is positive and varies
multiplicatively with width, depth, and unroll length, the three
coordinates that dominate it. Cold start uses the user's declared cost
hint if present and a uniform prior otherwise, with the cheap-first
initialization of CArBO carrying the first rounds until the model has
data. Failed and killed trials enter as censored observations at their
elapsed time (a lower bound on cost, not the cost), because dropping
them would teach the model that the expensive corner is cheap. The
model refits with the performance surrogate on the same cadence, and
its calibration is logged as a standing diagnostic beside the rung
correlation, so a badly-fit cost model is visible rather than silently
distorting proposals.

\section{Tier 2: the population line for large-parallelism RL}

Ray already ships PBT and a continuous-only PB2, so the work starts one
rung up: extend the explore step to mixed spaces (PB2-Mix's bandit over
categories, or directly the Casmopolitan-style kernels), then build
BG-PBT proper: the mixed-space trust-region searcher with ordinal-aware
kernels, generational restarts, and distillation hooks left to the
training code (Chapter~\ref{sec:population}). The official BG-PBT code
is available to crib from, and RLlib checkpointing supplies the surgery
PBT needs. This is the flagship for the large-parallelism regime, and
the product should encode the budget boundary from
\citet{eimer2023hyperparameters}: below a few dozen full-run
equivalents, steer users to DEHB or BO plus ASHA instead. Encoding the
boundary is not a footnote; it is the difference between a product that
embodies the evidence of Chapter~\ref{sec:population} and one that
merely implements its algorithms.

\section{Tier 3: the frontier, selectively}

Three bets with good expected value once Tier 1 is earning. A
learning-curve scheduler in the ifBO style, using the published
pre-trained freeze-thaw surrogate, exploiting our pause-and-resume
plumbing (Chapter~\ref{sec:multifidelity}); the decision problem it
solves is Figure~\ref{fig:freezethaw}, and the pre-trained surrogate
(Chapter~\ref{sec:transfer}) is why the bet is affordable. DEHB, which
the RL evidence recommends and which also gives us a non-GP engine for
awkward spaces; the upstream package is frozen
(Chapter~\ref{sec:libraries}), so this is an implementation against our
interfaces with the original as a reference. And HEBO-style robustness
(input and output warping, noise-aware acquisition ensembles) folded
into our GP searcher \citep{cowenrivers2022hebo}, which
Section~\ref{sec:gp-practical} argued is where the challenge-winning
gains actually came from. A CARBS-style cost-performance frontier is the
natural fourth once a scaling-bound project appears
(Chapter~\ref{sec:transfer}).

\section{What we deliberately skip}

Differentiable NAS and a bespoke zero-cost-proxy subsystem
(Chapter~\ref{sec:workloads}); gradient-based HPO through the training
loop; a meta-learned transformer surrogate trained by us (we adopt
pre-trained ones where they ship); and any LLM-driven configuration
proposer, which as of this writing has no controlled evidence of beating
the methods above at matched cost on our kinds of workloads. Each skip
is revisitable when the evidence changes; the skips are recorded here
precisely so that revisiting them is a decision rather than a drift.

\section{The calendar, and its dependencies}
\label{sec:calendar}

The tiers above say what and why; a build needs when and after what.
The calendar below assumes one dedicated engineer with part-time
research support for the campaign runs, an assumption the manager
owns, not the survey; with a second engineer the two independent
tracks (the multi-objective stack and the population line) run in
parallel and the total compresses by roughly a third. Weeks are
planning estimates in the same spirit as the cost column of
Table~\ref{tab:decisions}.

Weeks 1--2: freeze. The interfaces and store schema of
Chapter~\ref{sec:contracts} are frozen against a written review, the
pinned tasks of Chapter~\ref{sec:validation} are sized to their
compute envelope, and the layer-one harness (benchmark parity and
curve replay) is stood up first, because every later gate runs
through it.

Weeks 3--6: tier 0. The wrapped searchers and schedulers go behind
the interfaces; the run store lands with lineage and the standing
diagnostics; the Ray and local adapters pass checkpoint surgery
through. Gate: V01--V05 and V14 green. V04 tests the tier-0 GP and
TPE at this gate; tier-1 searchers face it at the tier-1 gate. The
tier-0 gate is deliberately mostly wiring tests plus the one method
floor check, because tier 0 makes no novelty claims, only composition
claims.

Weeks 7--12: tier 1. The multi-objective stack first (qLogNEHVI,
Chebyshev fallback, MO-ASHA with veto gates, weeks 7--9), then
priors, warm start, and cost-aware acquisition (weeks 10--11), then
the tier-1 campaign week: V04, V06, V09, V10, V11, V13, with V12
maturing later as the store fills. Dependencies bind the order:
MO-ASHA extends the tier-0 ASHA; PriorBand consumes its rungs; warm
start queries the store that tier 0 built; EI-per-cost decorates the
tier-0 GP searcher.

Weeks 13--20: the population line's prerequisite. The trust-region
mixed-space searcher is built once and serves twice, as the
high-dimensional single-objective searcher of
Chapter~\ref{sec:bayesian} and as the explore engine of the
population line; PB2-Mix lands on it in weeks 17--20 using Ray's
existing PBT machinery through the exploit and explore verbs of the
scheduler interface.

Weeks 21--30: the flagship. BG-PBT proper (generational restarts,
architecture coordinates, distillation hooks left to training code),
closing with the V08 boundary campaign around week 30, whose
demotion rule was fixed in Chapter~\ref{sec:population} long before
this code exists.

Weeks 31 onward: tier 3, each item behind its own gate. DEHB against
our interfaces, raced in V07; the ifBO scheduler only after V15
shows the bet live on public suites and only once pause-and-resume
is proven by the population work; HEBO-style robustness folded into
the GP searcher last, because it is an upgrade to a working thing
rather than a missing thing.

\section{The bill}

Engineering: roughly six weeks to the tier-0 gate, six more to the
tier-1 gate, eighteen for the population line through its campaign,
and about twelve of selective tier-3 work; call it forty-two
engineer-weeks, ten months of one engineer or six of two, before
tier-2 electives. These are the table's per-item costs rolled up,
not new estimates.

Compute: layer one is CPU-minutes and free in practice; each
layer-two tier gate is budgeted at the one-GPU-week class, and the
two priced campaigns (V07, and V08 with its both-sides-of-the-
boundary design) are of the same order each, per the sizing
arithmetic of Chapter~\ref{sec:validation}. The planning envelope
for the whole program is therefore a few GPU-weeks, spent in named
campaigns with recorded manifests rather than ad hoc runs.

Maintenance: the standing cost of the architecture is the one
Chapter~\ref{sec:architecture} priced, two small interfaces kept
current against Ray, plus pinned dependency versions reviewed on
their release cadence.

\section{Risks, and the tripwires that catch them}

Ray or Tune drifts under us. Tripwire: the layer-one replay and
parity tests run against every Ray upgrade before it is adopted;
the adapter boundary is the insurance, and its annual premium was
priced in Chapter~\ref{sec:architecture}.

Upstream churn in the searcher libraries. The audit of
Chapter~\ref{sec:libraries} already caught BoTorch deprecating
acquisition names and Optuna changing default samplers at 5.0;
mitigation is pinned versions, thin wrappers, and the store
recording the wrapped library's version in every study row, so a
result is never orphaned by an upgrade. Tripwire: V01 parity on any
dependency bump.

The population evidence fails to transfer to our cluster. That is
not a risk to manage but a measurement to take: V08 exists to take
it, and its failure path (demote the flagship, keep the rung
methods) was committed to in writing before the build began.

The seed protocol proves too expensive as a default. Layer three
records every opt-out; a high opt-out rate on real studies is the
tripwire, and the response is a cheaper default protocol argued in
an update to Chapter~\ref{sec:workloads}, not a silent loosening.

The pinned tasks are mis-sized. Week 1--2 sizing exists for this,
and Chapter~\ref{sec:validation}'s own overturn clause governs the
repair.

\section{How we will know the tuner works}

The full answer is now Chapter~\ref{sec:validation}: the register of
tests, the statistical protocol, the priced campaigns, and the rule
that a tier is done when its gate entries are green. What this
chapter adds is only the schedule on which those gates fire, weeks
6, 12, and 30 on the one-engineer calendar, and the standing rule
that no tier is declared done by the feeling of momentum.
